\subsection{Signature Factorization}

Identical parameter sequences and declared-type bundles appearing in multiple function signatures are a recurring source of drift and accidental redundancy. Viewed through the zero-error lens of Paper~1 and the derivation/SSOT lens of Paper~2, this is the same structural phenomenon we treat elsewhere: repeated observable semantics instantiate a representation fiber whose multiplicity quantifies the residual identification burden.

Concretely, let \(\pi:\mathrm{Sig}\to\mathrm{Sem}\) be the map that projects a declaration to its observable semantics (parameter-type sequence, return type, and declared behavioral contract). The fiber \(\pi^{-1}(s)\) collects all signatures indistinguishable to callers under the admissibility contract (callers observe only the declared signature + contract). Paper~1 shows that zero-error identification within a fiber requires auxiliary information of size at least \(\log_2|\pi^{-1}(s)|\) bits; Paper~2 shows that making the repeated signatures derived from a single declared interface (an abstract base class or canonical ABC) enforces coherence and removes the auxiliary burden.

Engineering corollary. If the same semantic parameter sequence appears in multiple places, factor it into a single abstract interface and make the others derived (inherit, delegate, or generate). Doing so is not merely style: it is a zero-error coherence transformation that removes an information-theoretic burden from the system and prevents accidental divergence between semantically identical call sites.

Tooling suggestion. A lightweight linter may enumerate candidate fibers by hashing parameter-type sequences and presenting groups with |\pi^{-1}(s)|>1 for developer review. This scan is conservative under finite-precision type representations and is an automatic instantiation of the Paper~1 counting law.

Template Method payoff. When factoring signatures into an abstract base class (ABC), move the shared pre/post logic into a single concrete method on the ABC and expose only the minimal varying behavior as abstract hook(s) that concrete classes override. This is the Template Method pattern: it concentrates behavioral degrees of freedom into small hooks, so derivation/SSOT arguments (Paper~2) become straightforward to apply and the information burden quantified in Paper~1 drops because callers observe a single canonical interface.

Minimal example (before / after): the `process` method is the concrete template; subclasses override `_handle` as the hook that varies. Shared sanitization and logging live once in the ABC, so changes affect all implementations at once.

\begin{verbatim}
# before: duplicated pre/post logic in each concrete class
class AProcessor:
    def process(self, data: dict) -> str:
        clean = sanitize(data)
        out = do_specific(clean)
        log_result(out)
        return out

class BProcessor:
    def process(self, data: dict) -> str:
        clean = sanitize(data)
        out = do_other(clean)
        log_result(out)
        return out

# after: Template Method. shared logic moved into ABC; varying step is a hook
from abc import ABC, abstractmethod

class Processor(ABC):
    def process(self, data: dict) -> str:            # concrete, caller-visible
        clean = self._sanitize(data)                # common pre
        out = self._handle(clean)                   # hook — override
        self._log_result(out)                       # common post
        return out

    def _sanitize(self, data: dict) -> dict:
        ...

    def _log_result(self, out: str) -> None:
        ...

    @abstractmethod
    def _handle(self, clean: dict) -> str:
        pass

class AProcessor(Processor):
    def _handle(self, clean: dict) -> str:
        return do_specific(clean)

class BProcessor(Processor):
    def _handle(self, clean: dict) -> str:
        return do_other(clean)
\end{verbatim}

\leanmeta{\LH{INF1}, \LH{DER2}}
